\subsection{Discussion}
\label{sec:discussion}
The grid architecture makes every stage cheap, training-free and
deterministic, which is what allowed exact cross-implementation
agreement and an offline replay of every count. Increasing the resolution
slightly raises detection recall and reduces merging of neighbouring
vehicles, but it also splits vehicles into fragments, which produce extra
tracks and crossing events. Classification with fixed thresholds
therefore favours fine grids, whereas counting favours coarse to medium
ones; after re-tuning, the classification advantage disappears and the
counting cost remains. The geometric features struggle with vans and buses
because the normalized image footprint of a vehicle depends on its
distance from the camera, on occlusion and on fragmentation as much as on
its class, and because vans (7\%) and buses (6\%) are rare among matched
vehicles. Physical dimensions from the calibration, perspective
correction, temporal shape statistics, appearance features and better
association (e.g.\ Kalman filtering with Hungarian matching
\cite{bewley2016,wojke2017}) are the natural next steps; none was
implemented here. The experiments establish the equivalence of the two
implementations on 12 videos, the slowdown from threading on one machine
for this code, and the resolution trade-offs on UA-DETRAC with fixed
ROIs. They do not establish generalization to other cameras or datasets,
the accuracy of the speed estimates, or counting accuracy sufficient for
deployment (test event-F1 $\le0.34$).

\subsection{Threats to Validity}
\label{sec:threats}
All thresholds were tuned for UA-DETRAC geometry, and the two ROIs are
fixed for every video, including intersections in the test set;
illumination, shadows and occlusion remain failure modes
\cite{abbas2021,chen2020}. Camera sites were inferred visually because
UA-DETRAC has no site labels; several merges are uncertain (listed in the
supplementary mapping), and the unseen-site test set still comes from the
same dataset. The weights of Eq.~\eqref{eq:objective} reflect our
judgment and were not optimized. The counter counts only downward
crossings. Speed estimation was not evaluated, uses a single scale without
perspective correction, and produced implausible values in a qualitative
run (Section~\ref{sec:speed}). Benchmark timings come from one laptop and
one run per workload and video; the reported spreads are between-video
variation, and run-to-run variation is visible (sequential loop
$8.03\pm0.29$\,ms in an earlier run vs.\ $9.41\pm1.84$\,ms here), although
the threads-only result has the same sign on all 12 videos. Repeated runs
\cite{kalibera2013} and other machines are needed. The grid study ran on
a different machine and OpenCV version. Tracks younger than 5 frames are
never scored in the classification evaluation.

